\section{Commutativity in the World of Qubit Operators}

\subsection{Pinching a Matrix by a Hermitian Qubit Operator}

\begin{theorem}[Pinching any matrix by a Hermitian Matrix with 2 eigenvalues i.e. linear combination of a projector and its complement]\label{thm:pinch}
    Let $M = \begin{pmatrix}
        M_{1,1} & M_{1,2}\\
        M_{2,1} & M_{2,2}
    \end{pmatrix}$ be any $2d \times 2d$ matrix and let $N$ be a $2d \times 2d$ Hermitian matrix with eigenvalues $\lambda_1 > \lambda_2$. Write
    \[ N= \lambda_1 P + \lambda_2 (I - P) = U \begin{pmatrix}
        \lambda_1\: I_{d \times d} & 0\\ 0 & \lambda_2\: I_{d \times d}
    \end{pmatrix} U^\dagger , \] where $P$ is the projector into the $\lambda_1$-eigenspace and $U$ is the unitary diagonalizing $N$. Define:
    \begin{equation}
        M' := P\ M\ P + (I - P)\ M\ (I - P)
    \end{equation}
    Then, $[M', N] = 0$ and $\|M' - M \| = \frac{1}{\lambda_1 - \lambda_2} \|[M, N]\|$.
\end{theorem}

\begin{proof}
Without loss of generality, assume $U = I$, i.e. assume we have changed basis so that $N$ is diagonal. Let us explicitly calculate the commutator of $M$ with $N$.
\begin{align*}
[M, N] &= M\ N - N\ M\\
&= \begin{pmatrix} \lambda_1 M_{1,1} - \lambda_1 M_{1,1} & \lambda_2 M_{1,2} - \lambda_1 M_{1,2} \\ \lambda_1 M_{2,1} - \lambda_2 M_{2,1} & \lambda_2 M_{2,2} - \lambda_2 M_{2,2}
\end{pmatrix}\\
&= \begin{pmatrix} 0 & - \left(\lambda_1 - \lambda_2\right) M_{1,2} \\ \left(\lambda_1 - \lambda_2\right) M_{2,1} & 0
\end{pmatrix}\\
&= \left(\lambda_1 - \lambda_2\right)\begin{pmatrix} 0 & - M_{1,2} \\  M_{2,1} & 0
\end{pmatrix}\\
\end{align*}

Using this, we will now prove the two conclusions of the Theorem. First, since $M'$ is block-diagonal by construction, it follows from the calculation above applied to $M'$ instead of $M$ that $[M', N] = 0$. Now, to bound the rounding error $\|M' - M\|$, we calculate it explicitly.
\begin{align*}
\|M - M' \| &= \| (M - M')(2P - I)\|\\
&= \left\|\begin{pmatrix}
            0 & -M_{1,2}\\
            M_{2,1} & 0
        \end{pmatrix} \right\|\\
&= \frac{\lambda_1 - \lambda_2}{\lambda_1 - \lambda_2} \left\|\begin{pmatrix}
            0 & -M_{1,2}\\
            M_{2,1} & 0
        \end{pmatrix} \right\|\\  
&= \frac{1}{\lambda_1 - \lambda_2}\left\|(\lambda_1 - \lambda_2)\begin{pmatrix}
            0 & -M_{1,2}\\
            M_{2,1} & 0
        \end{pmatrix} \right\|\\
&= \frac{\|[M, N]\|}{\lambda_1 - \lambda_2}
\end{align*}

\end{proof}

\subsection{Snapping a Hermitian Matrix into a Multiple of the Identity}
\begin{theorem}[Globally snapping a hermitian matrix with 2 eigenvalues into identity]
\label{lemma:snapping-hermitian-qubit-hermitian-operator}
Let $B$ be a Hermitian matrix of the form $B = \lambda_1 P + \lambda_2 (I - P) = U \begin{pmatrix}
        \lambda_1\: I_{d \times d} & 0\\ 0 & \lambda_2\: I_{d \times d}
    \end{pmatrix} U^\dagger$, where $P$ is an orthogonal projector and $U$ is a unitary. Let
    $B' = \lambda_1\: I_{2d \times 2d}$.
Then, $\|B' - B \| = |\lambda_1 - \lambda_2|$.
\end{theorem}

\begin{proof}
\begin{align*}
\|B' - B \| &= \left\|\begin{pmatrix}
        0 & 0\\ 0 & \left(\lambda_1 - \lambda_2\right)\: I_{d \times d}
    \end{pmatrix}\right\| = |\lambda_1 - \lambda_2|
\end{align*}
\end{proof}
\subsection{Transitivity of (Almost) Commutativity of Qubit Operators}

\begin{lemma}[Gapped spectra imply transitivity of commutativity]\label{lem:transitivity}
Let $A,B,C \in \mathrm{Herm}(\mathbb{C}^{2 \times 2})$ and suppose that $B$ has non-zero spectral gap $\Delta(B)=|\lambda_1(B) - \lambda_2(B)| > 0$. Then, 
$$
[A,B] = 0 \quad \text{ and } \quad [B,C] = 0 \quad \implies \quad [A,C]=0.
$$
\end{lemma}
\begin{proof}
Firstly, since $\Delta(B) > 0$, $B$ has a unique eigenbasis in which it is diagonal. Let $\{P, P^\perp\}$ be the orthogonal projectors onto the two elements of this basis. 
According to Fact~\ref{fact:sim-diag}, since $A$ commutes with $B$, and $C$ commutes with $B$, and since~\footnote{Without that condition, we can only conclude that $A$ and $B$ can be simultaneously diagonalized and that $B$ and $C$ can be simultaneously diagonalized, but not necessarily all three can be simultaneously diagonalized. If $B$ is a multiple of the identity, then any two non-commuting matrices will be a counterexample.} $\Delta(B) > 0$, we can diagonalize all the three matrices in the unique basis $\{P, P^\perp\}$ as follows:
\begin{equation}
    A = a_1 P + a_2 P^\bot
\end{equation}
\begin{equation}
    B = b_1 P + b_2 P^\bot
\end{equation}
\begin{equation}
    C = c_1 P + c_2 P^\bot
\end{equation}
Thus, we can see that:
\begin{align*}
    [A, C] &= \left[a_1 P + a_2 P^\bot,  c_1 P + c_2 P^\bot\right]\\
    &= a_1c_1 [P, P] + a_1 c_2 [P, P^\bot] + a_2 c_1 [P^\bot, P] + a_2 c_2 [P^\bot, P^\bot]\\
    &= 0.
\end{align*}
\end{proof}